\subsection{Posets}
\begin{frame}{Partial ordering}
    Ligner på ekvivalensrelasjoner, men er anti-symmetrisk istedenfor symmetrisk\\
    Mengder med en partial ordering relasjon på seg kalles for \emph{posets}\\

    \begin{columns}
        \begin{column}{0.4\textwidth}
            En ekvivalenserelasjon:
            \begin{tikzcd}[ampersand replacement=\&]
                a \arrow[d, bend right] \arrow[loop, distance=2em, in=125, out=55] \arrow[r, bend left]
                \& b \arrow[loop, distance=2em, in=125, out=55] \arrow[ld, bend right] \arrow[l, bend left]
                \& c \arrow[loop, distance=2em, in=125, out=55] \\
                d \arrow[ru, bend right] \arrow[u, bend right] \arrow[loop, distance=2em, in=305, out=235]
                \& e \arrow[r, bend right] \arrow[loop, distance=2em, in=305, out=235]
                \& f \arrow[loop, distance=2em, in=305, out=235] \arrow[l, bend right]
            \end{tikzcd}
        \end{column}
        \begin{column}{0.4\textwidth}
            En partial order relasjon:
            \begin{tikzcd}[ampersand replacement=\&]
                a \arrow[d] \arrow[loop, distance=2em, in=125, out=55] \arrow[r]
                \& b \arrow[loop, distance=2em, in=125, out=55] \arrow[ld]
                \& c \arrow[loop, distance=2em, in=125, out=55] \\
                d \arrow[loop, distance=2em, in=305, out=235]
                \& e \arrow[r] \arrow[loop, distance=2em, in=305, out=235]
                \& f \arrow[loop, distance=2em, in=305, out=235] 
            \end{tikzcd}
        \end{column}
    \end{columns}
    To elementer kalles for \emph{sammenlignbar} hvis de har en pil mellom seg.\\
    Hvis alle to elementer kan sammenlignes, snakker vi om en \emph{total ordering}
\end{frame}